\section{Conclusion and Future Work}\label{sec:conclusion}

In conclusion, the introduction of \ours marks a significant step forward in the development of autonomous digital agents, addressing critical gaps in existing interactive learning environments. 
By providing a rich, realistic setting that spans multiple operating systems, interfaces, and applications, \ours not only broadens the scope of tasks digital agents can perform but also enhances their potential for real-world application. 
Despite the promise shown by advancements in vision-language models, evaluations within \ours reveal notable challenges in agents' abilities, particularly in GUI understanding and operational knowledge, pointing to essential areas for future research and development.

We identify several potential directions for community development and progress toward general-purpose agents for computer operation:
\paragraph{Enhancing VLM capabilities for efficient and robust GUI interactions}
For foundation model development, we need to boost the efficiency of our models, enabling them to process much longer contexts and perform inference computations efficiently, akin to the robotics community~\citep{brohan2022rt, brohan2023rt} to better handle real-world cases.
Enhancements in VLMs' GUI grounding capabilities that is robust to application windows changes and  are also sought, focusing on the accurate understanding and generation of precise actions aligned with given instructions. 
Moreover, amplifying VLMs' ability to comprehend context in the form of images is a pivotal goal, since it is crucial to enable history encoding using images so that we can build memory and reflection upon that.
These improvements may require more efforts in the upstream pre-training stage, downstream fine-tuning stage, and even in the model structure itself, as pointed out in previous work~\citep{deng2023mind2web, hong2023cogagent, lu2024weblinx}.

\paragraph{Advancing agent methodologies for exploration, memory, and reflection}
The next-level approach encompasses designing more effective agent architectures that
augment the agents' abilities to explore autonomously and synthesize their findings. 
The agents face challenges in leveraging lengthy raw observation and action records. It's fascinating to explore novel methods for encoding this history, incorporating efficient memory and reflection solutions to condense contextual information and aid the agent in extracting key information. Additionally, integrating knowledge grounding into (V)LLM agents through memory mechanisms is a promising avenue as well.
Moreover, practice GUI assistants also require features of personalization and customization. These features rely on techniques such as user profiling and retaining memories from long-term user-assistant interactions.
Additionally, crafting protocols specifically for digital agents operating within GUI and CLI interfaces aims at facilitating efficient actions is also an essential thing for the feasibility of general-purpose digital agents in the mid-short term.

\paragraph{Addressing the safety challenges of agents in realistic environments}
The safety of agents is a critical issue if applying a built agent in fully realistic environments, the developed universal digital agent could potentially be used to bypass CAPTCHA systems in the future, as noted in \citep{searles2023empirical}. 
However, due to the currently limited capabilities of agents, we have not observed any harmful and damaging behaviors during our experiments, an automatic agent has the opportunity to damage patent rights, abuse accounts, attempt to exploit software vulnerabilities to create viruses, or engage in attacks. 
Currently, we adopt virtual machines to make it difficult for developing digital agents to cause irreversible damage to our host machines. 
However, there still lacks a reliable metric to assess the safety of an agent developed in an isolated environment. 
The current evaluation functions mainly focus on the results closely regarding the task instructions, assess only the correctness of task completion, and pay little attention to potential unnecessary damaging actions of agents.
Owing to the complexity of a complete computer environment, we didn't work out an efficient way to detect the latent side effects of the agent.
Consequently, how to assess and control potential behaviors in open and real environments through environmental constraints and agent training is an important further direction of research.

\paragraph{Expanding and refining data and environments for agent development}
In terms of datasets and environments, we can broaden the scope to cover more specialized domains, including real-sector needs in healthcare, education, industry, transportation, and personalized requirements. 
Efforts can be made to ensure our environment's seamless deployment across various hardware and software settings. 
The variance of a11y tree quality across different applications is also noticed. Although
the problem is not remarkable in the applications currently included, there is no guarantee
of that the application developers obey the a11y convention and offer clear and meaningful 
descriptions for GUI elements. More intelligent approaches to filter redundant a11y tree elements
and to handle latently missing elements deserve careful investigation as well.
We also highlight the necessity of a painless data collection method, allowing for the effortless acquisition of computer operation data and its transformation into agent capabilities.
